\documentclass[conference]{IEEEtran}
\IEEEoverridecommandlockouts

\usepackage{cite}
\usepackage{amsmath,amssymb,amsfonts}
\usepackage{algorithmic}
\usepackage{graphicx}
\usepackage{textcomp}
\usepackage{xcolor}
\usepackage{hyperref}

\def\BibTeX{{\rm B\kern-.05em{\sc i\kern-.025em b}\kern-.08em
    T\kern-.1667em\lower.7ex\hbox{E}\kern-.125emX}}

\begin{document}

\title{Beyond ReAct: Evaluating and Improving LLM Tool Agents on Real-World Benchmarks}

\author{
    \IEEEauthorblockN{Muntaha Asim}
    \IEEEauthorblockA{muntahaasim79@gmail.com}
    \and
    \IEEEauthorblockN{Ammar Kashif}
    \IEEEauthorblockA{ammarkashif00@gmail.com}
}

\maketitle

\begin{abstract}
The automation of machine learning (ML) experimentation has long been a goal
of the artificial intelligence research community. Recent advances in large
language models (LLMs) have opened new possibilities for building autonomous
agents capable of performing complex, multi-step tasks such as writing and
executing code, analyzing results, and iterating toward improved model
performance. This paper reviews the current state of LLM-based agentic systems
applied to ML experimentation, with a focus on benchmarking, agent design, and
key limitations. We examine foundational frameworks such as ReAct and
Reflexion, survey benchmark efforts including MLAgentBench and AgentBench, and
analyze the performance of agents powered by models such as GPT-4, Claude, and
Gemini across diverse ML tasks. Despite promising results, significant
challenges remain, including hallucination, poor long-term planning, and limited
generalization to tasks unseen during model pretraining. This review identifies
these open problems and outlines directions for future research toward more
reliable and capable ML experimentation agents.
\end{abstract}

\begin{IEEEkeywords}
large language models, autonomous agents, machine learning experimentation,
benchmarking, ReAct, MLAgentBench, AutoML
\end{IEEEkeywords}


%-------------------------------------------------------
\section{Introduction}
%-------------------------------------------------------

Machine learning research is fundamentally driven by experimentation. A typical
ML workflow involves formulating a hypothesis, selecting an appropriate model
architecture, implementing and executing training code, evaluating performance
metrics, and iterating based on the results. This process demands not only deep
technical knowledge but also the ability to reason across long time horizons,
debug code, and draw meaningful conclusions from empirical observations.
Traditionally, this workflow has been carried out entirely by human researchers,
making it both time-intensive and expertise-dependent.

The emergence of large language models (LLMs) such as GPT-4 \cite{openai2023gpt4},
Claude \cite{anthropic2023claude}, and Gemini \cite{anil2023gemini} has introduced
a fundamentally new paradigm in AI. These models exhibit remarkable capabilities
in natural language understanding, code generation, and structured reasoning.
Building on these capabilities, recent research has explored the possibility of
constructing autonomous agents --- systems that can perceive their environment,
plan sequences of actions, and execute those actions with minimal human
intervention \cite{wang2023voyager,park2023generative}.

LLM-based agents represent a promising avenue for automating ML experimentation
end-to-end. Rather than relying on narrowly scoped AutoML pipelines that optimize
a fixed set of hyperparameters, agentic systems can interact with a file system,
write and modify training scripts, execute code, and adaptively revise their
strategies based on observed outcomes. This flexibility allows them to operate
in open-ended settings that more closely resemble real-world research workflows.

However, existing agentic systems face significant limitations when applied to
complex ML tasks. Key challenges include the tendency to hallucinate results ---
claiming performance improvements without empirical verification --- poor planning
over extended interaction sequences, and weak generalization to problem domains
that were not represented in the model's pretraining data \cite{huang2024mlagentbench}.
These limitations reduce the reliability of LLM-based agents in settings where
correctness and reproducibility are critical.

This paper reviews the current landscape of LLM-based agents for ML
experimentation. We examine the design of agent frameworks, the benchmarks used
to evaluate them, and the empirical performance reported across a range of tasks
and models. By synthesizing findings from recent literature, we aim to
characterize the state of the field, identify persistent gaps, and motivate
future research directions.


%-------------------------------------------------------
\section{Literature Review}
%-------------------------------------------------------

\subsection{LLM-Based Agent Frameworks}

The development of general-purpose LLM agents has been an active area of
research. Yao et al. introduced ReAct \cite{yao2023react}, a framework that
interleaves reasoning traces with action execution, enabling agents to maintain
a coherent internal state while interacting with external environments. ReAct
demonstrated that prompting LLMs to reason before acting significantly improves
task completion across question answering and interactive decision-making
benchmarks. This design principle has since been widely adopted in subsequent
agentic systems.

Shinn et al. proposed Reflexion \cite{shinn2023reflexion}, which extends ReAct
by introducing verbal reinforcement learning. In Reflexion, agents generate
natural language reflections on their previous failures and use these reflections
to improve future behavior without updating model weights. This approach proved
particularly effective in tasks requiring iterative refinement, making it
directly relevant to ML experimentation workflows where agents must learn from
failed experiments.

Park et al. explored generative agents \cite{park2023generative} that simulate
believable human behavior through memory, reflection, and planning modules.
While not directly applied to ML experimentation, their architecture demonstrated
that LLMs can maintain coherent long-term behavior across extended interaction
sequences --- a key capability required for autonomous research workflows.

\subsection{Benchmarking Agentic Capabilities}

Evaluating the capabilities of LLM-based agents requires carefully constructed
benchmarks that capture the complexity of real-world tasks. Liu et al. proposed
AgentBench \cite{liu2024agentbench}, a comprehensive evaluation framework that
tests LLMs across eight distinct environments including web browsing, database
interaction, and operating system tasks. AgentBench revealed substantial
performance gaps between different model families and highlighted the difficulty
of sustaining consistent reasoning across long task horizons.

Zhou et al. introduced WebArena \cite{zhou2023webarena}, a realistic web-based
environment for evaluating agents on tasks that require navigation, form
submission, and information retrieval across live web applications. Their
findings showed that even state-of-the-art LLMs achieve relatively low success
rates on realistic web tasks, suggesting that agentic performance degrades
significantly as task complexity increases.

Kinniment et al. proposed ARA \cite{kinniment2023ara}, a benchmark targeting
high-stakes autonomous tasks that require extended planning and execution. ARA
emphasizes the importance of safety and reliability in agentic systems,
particularly when agents are given access to real computational resources.

\subsection{Agents for Machine Learning Experimentation}

Huang et al. introduced MLAgentBench \cite{huang2024mlagentbench}, the first
benchmark specifically designed to evaluate LLM-based agents on ML
experimentation tasks. MLAgentBench comprises 13 tasks spanning diverse domains
including image classification, natural language processing, graph learning, and
code optimization. For each task, an agent is provided with starter code,
dataset files, and a task description, and must iteratively improve model
performance through file editing and code execution.

Experiments in MLAgentBench evaluated agents powered by GPT-4, GPT-4-turbo,
Claude v1.0, Claude v2.1, Claude v3 Opus, Gemini Pro, and Mixtral. Claude v3
Opus achieved the highest average success rate of 37.5\%, while other models
performed considerably lower. Importantly, success rates dropped sharply on
tasks introduced after the models' pretraining cutoff dates, with some Kaggle
challenges yielding 0\% success across all models \cite{huang2024mlagentbench}.
This finding highlights a critical limitation: agents relying solely on
pretrained knowledge struggle to generalize to genuinely novel problem settings.

The agent design in MLAgentBench incorporates structured reasoning through four
components: \textit{Reflection}, \textit{Research Plan and Status},
\textit{Fact Check}, and \textit{Thought}. The Fact Check component was
introduced specifically to mitigate hallucination --- a prevalent failure mode
in which agents falsely claim performance improvements without executing the
relevant code. While this design reduced hallucination frequency, it did not
eliminate it entirely, indicating that structural prompting alone is insufficient
for reliable agentic behavior.

\subsection{AutoML and LLM-Assisted Optimization}

Prior to the emergence of fully agentic systems, the ML community explored
AutoML as a means of reducing human effort in model development. Neural
Architecture Search (NAS) \cite{elsken2019nas} automates the design of neural
network architectures, while hyperparameter optimization frameworks such as
AutoML \cite{he2021automl} search over predefined configuration spaces. These
approaches, however, are constrained to fixed search spaces and lack the
flexibility to reason about arbitrary code or adapt to novel problem formulations.

Zhang et al. proposed AutoML-GPT \cite{zhang2023automlgpt}, which prompts LLMs
with data and model cards to predict training outcomes and guide hyperparameter
selection. While demonstrating the utility of LLMs for structured optimization
tasks, AutoML-GPT operates within a narrower scope compared to fully interactive
agentic systems. Similarly, MLcopilot \cite{zhang2023mlcopilot} uses LLMs to
recommend categorized hyperparameter settings based on accumulated past
experience, but does not support open-ended code interaction. These approaches
collectively suggest a continuum from constrained AutoML pipelines toward fully
autonomous agentic systems, with LLM-based agents representing the open-ended
end of this continuum.

\subsection{Identified Gaps and Open Problems}

Despite notable progress, several important gaps remain in the literature.
First, existing benchmarks primarily evaluate task success rates without deeply
analyzing the quality of the agent's reasoning process. Second, most current
work focuses on single-agent architectures, leaving multi-agent collaboration
for ML experimentation largely unexplored. Third, the impact of context window
limitations on long-horizon task performance has not been systematically studied.
Finally, the safety implications of agents that autonomously write and execute
code on real systems remain underexplored, despite being directly relevant to
deployment in research environments. Addressing these gaps represents a
significant opportunity for future work in agentic AI.


%-------------------------------------------------------
\begin{thebibliography}{00}

\bibitem{openai2023gpt4}
OpenAI, ``GPT-4 technical report,'' \textit{arXiv}, abs/2303.08774, 2023.

\bibitem{anthropic2023claude}
Anthropic, ``Introducing Claude,'' 2023. [Online]. Available:
https://www.anthropic.com/index/introducing-claude

\bibitem{anil2023gemini}
G.~T.~G.~R. Anil et al., ``Gemini: A family of highly capable multimodal
models,'' \textit{arXiv}, abs/2312.11805, 2023.

\bibitem{wang2023voyager}
G.~Wang, Y.~Xie, Y.~Jiang, A.~Mandlekar, C.~Xiao, Y.~Zhu, L.~J. Fan, and
A.~Anandkumar, ``Voyager: An open-ended embodied agent with large language
models,'' \textit{arXiv}, abs/2305.16291, 2023.

\bibitem{park2023generative}
J.~S. Park, J.~O'Brien, C.~J. Cai, M.~R. Morris, P.~Liang, and
M.~S. Bernstein, ``Generative agents: Interactive simulacra of human
behavior,'' in \textit{Proc. 36th Annual ACM Symposium on User Interface
Software and Technology (UIST)}, 2023.

\bibitem{huang2024mlagentbench}
Q.~Huang, J.~Vora, P.~Liang, and J.~Leskovec, ``MLAgentBench: Evaluating
language agents on machine learning experimentation,'' in \textit{Proc. 41st
International Conference on Machine Learning (ICML)}, PMLR 235, 2024.

\bibitem{yao2023react}
S.~Yao, J.~Zhao, D.~Yu, N.~Du, I.~Shafran, K.~Narasimhan, and Y.~Cao,
``ReAct: Synergizing reasoning and acting in language models,'' in
\textit{International Conference on Learning Representations (ICLR)}, 2023.

\bibitem{shinn2023reflexion}
N.~Shinn, F.~Cassano, A.~Gopinath, K.~R. Narasimhan, and S.~Yao,
``Reflexion: Language agents with verbal reinforcement learning,'' in
\textit{Thirty-seventh Conference on Neural Information Processing
Systems (NeurIPS)}, 2023.

\bibitem{liu2024agentbench}
X.~Liu et al., ``AgentBench: Evaluating LLMs as agents,'' in \textit{The
Twelfth International Conference on Learning Representations (ICLR)}, 2024.

\bibitem{zhou2023webarena}
S.~Zhou et al., ``WebArena: A realistic web environment for building autonomous
agents,'' \textit{arXiv}, abs/2307.13854, 2023.

\bibitem{kinniment2023ara}
M.~Kinniment et al., ``Evaluating language-model agents on realistic autonomous
tasks,'' \textit{arXiv}, abs/2312.11671, 2023.

\bibitem{elsken2019nas}
T.~Elsken, J.~H. Metzen, and F.~Hutter, ``Neural architecture search: A
survey,'' \textit{Journal of Machine Learning Research}, vol.~20, no.~1,
pp.~1997--2017, 2019.

\bibitem{he2021automl}
X.~He, K.~Zhao, and X.~Chu, ``AutoML: A survey of the state-of-the-art,''
\textit{Knowledge-Based Systems}, vol.~212, p.~106622, 2021.

\bibitem{zhang2023automlgpt}
S.~Zhang, C.~Gong, L.~Wu, X.~Liu, and M.~Zhou, ``AutoML-GPT: Automatic machine
learning with GPT,'' \textit{arXiv}, abs/2305.02499, 2023.

\bibitem{zhang2023mlcopilot}
L.~Zhang, Y.~Zhang, K.~Ren, D.~Li, and Y.~Yang, ``MLcopilot: Unleashing the
power of large language models in solving machine learning tasks,''
\textit{arXiv}, abs/2304.14979, 2023.

\end{thebibliography}

\end{document}
